\setupcolors[state=start]
\setupcolor[hex]
\definecolor[darkblue][h=003366]
\definecolor[darkgray][h=e6e6e6]
\definecolor[redline][h=d9534f]
\definecolor[blueline][h=5bc0de]
\definecolor[greenline][h=5cb85c]

\setuphead[section][style=\bfa, color=darkblue, before=\blank, after=\blank]
\setuphead[subsection][style=\bf, color=black, before=\blank, after=\blank]
\setuppapersize[A4][A4]
\setupbodyfont[11pt]
\setuplayout[width=middle, height=middle, topspace=2cm, bottomspace=2cm]

\starttext
\title{\bf Modelo OPP v4: Asignación de lotes a pickers}
\blank[big]

\startsection[title={Introducción}]
El modelo OPP v4 es una formulación de programación lineal entera para el problema de \emph{Order Picking Problem} que asigna lotes de órdenes a pickers y modela la secuenciación de lotes en cada picker. 
Esta versión extiende modelos anteriores con:
\startitemize
    \item {\bf Asignación única:} Cada lote se asigna de manera exclusiva a un solo picker.
    \item {\bf Variables de ruta:} Representación detallada para el trayecto de cada lote en el almacén.
    \item {\bf Tiempo de ejecución:} Incorpora tiempos de setup, picking, viaje y servicio en los nodos.
    \item {\bf Earliness / Tardiness:} Propiedades y penalizaciones según los tiempos de entrega de las órdenes.
    \item {\bf Secuenciación estricta:} Garantía de no solapamiento entre los lotes asignados a un mismo picker.
\stopitemize
\stopsection

\startsection[title={Conjuntos}]
\startformula
\startalign
    \NC \mathcal{N} \NC = \text{Órdenes de clientes} \NR
    \NC \mathcal{V} \NC = \text{Locaciones del almacén, incluido } v_1 \text{ y } v_f \NR
    \NC \mathcal{P} \NC = \text{Pickers} \NR
    \NC \mathcal{B} \NC = \text{Lotes factibles de pedidos} \NR
    \NC \mathcal{E} \NC = \text{Arcos válidos del grafo de locaciones} \NR
\stopalign
\stopformula
\stopsection

\startsection[title={Parámetros}]
\startformula
\startalign
    \NC v_1 \NC = \text{Nodo depósito inicial} \NR
    \NC v_f \NC = \text{Nodo final} \NR
    \NC W_{full} \NC = \text{Peso del tiempo de completitud de los lotes} \NR
    \NC W_{ea\_ta} \NC = \text{Peso de la penalización total por adelanto y tardanza} \NR
    \NC W_{ea} \NC = \text{Peso del adelanto (gamma)} \NR
    \NC W_{ta} \NC = \text{Peso de la tardanza (delta)} \NR
    \NC B_{time} \NC = \text{Tiempo por unidad de distancia} \NR
    \NC B_{setup} \NC = \text{Tiempo de preparación por lote} \NR
    \NC B_{picking} \NC = \text{Tiempo de picking por unidad de capacidad} \NR
    \NC C_v \NC = \text{Capacidad del vehículo de picking} \NR
    \NC M \NC = \text{Constante Big-M suficientemente grande} \NR
    \NC C_{n} \NC = \text{Capacidad requerida por la orden } n \in \mathcal{N} \NR
    \NC Dd_{n} \NC = \text{Due date de la orden } n \in \mathcal{N} \NR
    \NC d_{ij} \NC = \text{Distancia entre locaciones } i,j \in \mathcal{V} \text{ con } (i,j) \in \mathcal{E} \NR
    \NC S_{i} \NC = \text{Tiempo de servicio en la locación } i \in \mathcal{V} \NR
    \NC A_{ni} \NC = \startmathcases
                    \NC 1 \MC \text{si la orden } n \in \mathcal{N} \text{ requiere visitar } i \in \mathcal{V} \NR
                    \NC 0 \MC \text{en otro caso} \NR
                 \stopmathcases \NR
    \NC Bn_{nb} \NC = \startmathcases
                    \NC 1 \MC \text{si la orden } n \in \mathcal{N} \text{ pertenece al lote } b \in \mathcal{B} \NR
                    \NC 0 \MC \text{en otro caso} \NR
                 \stopmathcases \NR
    \NC R_{b} \NC = \left| \left\{ i \in \mathcal{V} \setminus \{v_1,v_f\} : \sum_{n \in \mathcal{N}} A_{ni} \cdot Bn_{nb} \ge 1 \right\} \right| \NR
\stopalign
\stopformula
\stopsection

\startsection[title={Variables}]
\startformula
\startalign
    \NC ta_{n} \NC = \text{Tardanza de la orden } n \in \mathcal{N} \NR
    \NC ea_{n} \NC = \text{Adelanto de la orden } n \in \mathcal{N} \NR
    \NC t\_exec_{b} \NC = \text{Tiempo de ejecución del lote } b \in \mathcal{B} \NR
    \NC v_{b} \NC = \text{Tiempo de completitud del lote } b \in \mathcal{B} \NR
    \NC u_{pb} \NC = \text{Tiempo de procesamiento del lote } b \in \mathcal{B} \text{ por el picker } p \in \mathcal{P} \NR
    \NC v_{pb} \NC = \text{Tiempo de completitud del lote } b \in \mathcal{B} \text{ por el picker } p \in \mathcal{P} \NR
    \NC z_{ijb} \NC = \startmathcases
                    \NC 1 \MC \text{si el arco } (i,j) \in \mathcal{E} \text{ es usado en el lote } b \in \mathcal{B} \NR
                    \NC 0 \MC \text{en otro caso} \NR
                 \stopmathcases \NR
    \NC f_{ijb} \NC = \text{Flujo de unidades sobre el arco } (i,j) \in \mathcal{E} \text{ en el lote } b \in \mathcal{B} \NR
    \NC X_{pb} \NC = \startmathcases
                    \NC 1 \MC \text{si el lote } b \in \mathcal{B} \text{ es asignado al picker } p \in \mathcal{P} \NR
                    \NC 0 \MC \text{en otro caso} \NR
                 \stopmathcases \NR
    \NC Z_{bk} \NC = \startmathcases
                    \NC 1 \MC \text{variable auxiliar para pares } b < k \NR
                    \NC 0 \MC \text{en otro caso} \NR
                 \stopmathcases \NR
\stopalign
\stopformula
\stopsection

\startsection[title={Función objetivo}]
\placeformula
\startformula
\min \left( W_{full} \cdot \sum_{p \in \mathcal{P}} \sum_{b \in \mathcal{B}} v_{pb} + W_{ea\_ta} \cdot \left( W_{ea} \cdot \sum_{n \in \mathcal{N}} ea_{n} + W_{ta} \cdot \sum_{n \in \mathcal{N}} ta_{n} \right) \right)
\stopformula
\stopsection

\startsubsection[title={Rutas de lote}]
\startformula
\startalign[n=3, align={left,left,left}]
    \NC \sum_{j:(v_1,j) \in \mathcal{E}} z_{v_1 j b} = 1; \NC \quad \forall b \in \mathcal{B} \NC \quad \text{(Cada lote sale una vez del depósito)} \NR
    \NC \sum_{i:(i,v_f) \in \mathcal{E}} z_{i v_f b} = 1; \NC \quad \forall b \in \mathcal{B} \NC \quad \text{(Cada lote entra una vez al nodo final)} \NR
    \NC \sum_{i:(i,v_1) \in \mathcal{E}} z_{i v_1 b} = 0; \NC \quad \forall b \in \mathcal{B} \NC \quad \text{(No hay retorno al depósito inicial)} \NR
    \NC \sum_{j:(v_f,j) \in \mathcal{E}} z_{v_f j b} = 0; \NC \quad \forall b \in \mathcal{B} \NC \quad \text{(No hay salida desde el nodo final)} \NR
\stopalign
\stopformula
\stopsubsection

\startsubsection[title={Visita y conservación de flujo}]
\startformula
\startalign[n=3, align={left,left,left}]
    \NC \sum_{i:(i,j) \in \mathcal{E}} z_{i j b} \ge \min\left(1, \sum_{n \in \mathcal{N}} A_{n j} \cdot Bn_{n b}\right); \NC \quad \forall b \in \mathcal{B},\ j \in \mathcal{V} \setminus \{v_1,v_f\} \NC \quad \text{(Entrada a locaciones requeridas)} \NR
    \NC \sum_{i:(j,i) \in \mathcal{E}} z_{j i b} \ge \min\left(1, \sum_{n \in \mathcal{N}} A_{n j} \cdot Bn_{n b}\right); \NC \quad \forall b \in \mathcal{B},\ j \in \mathcal{V} \setminus \{v_1,v_f\} \NC \quad \text{(Salida de locaciones requeridas)} \NR
    \NC \sum_{g:(g,i) \in \mathcal{E}} z_{g i b} - \sum_{j:(i,j) \in \mathcal{E}} z_{i j b} = 0; \NC \quad \forall b \in \mathcal{B},\ i \in \mathcal{V} \setminus \{v_1,v_f\} \NC \quad \text{(Conservación en nodos intermedios)} \NR
    \NC \sum_{j:(j,i) \in \mathcal{E}} z_{j i b} \le 1; \NC \quad \forall b \in \mathcal{B},\ i \in \mathcal{V} \setminus \{v_1\} \NC \quad \text{(A lo sumo un arco de entrada por nodo)} \NR
    \NC \sum_{j:(i,j) \in \mathcal{E}} z_{i j b} \le 1; \NC \quad \forall b \in \mathcal{B},\ i \in \mathcal{V} \setminus \{v_f\} \NC \quad \text{(A lo sumo un arco de salida por nodo)} \NR
\stopalign
\stopformula
\stopsubsection

\startsubsection[title={Subtour elimination y flujo de carga}]
\startformula
\startalign[n=3, align={left,left,left}]
    \NC f_{i j b} - R_{b} \cdot z_{i j b} \le 0; \NC \quad \forall b \in \mathcal{B},\ (i,j) \in \mathcal{E} \NC \quad \text{(Límite MTZ subtours)} \NR
    \NC \sum_{j:\ (v_1,j) \in \mathcal{E}} f_{v_1 j b} = R_{b}; \NC \quad \forall b \in \mathcal{B} \NC \quad \text{(Flujo total desde el depósito)} \NR
    \NC \sum_{ {g \atop (g,i) \in \mathcal{E}} } f_{g i b} - \sum_{ {j \atop (i,j) \in \mathcal{E}} } f_{i j b} = \min\left(1, \sum_{n \in \mathcal{N}} A_{n i} \cdot Bn_{n b}\right); \NC \quad \forall b \in \mathcal{B},\ i \in \mathcal{V} \setminus \{v_1,v_f\} \NC \quad \text{(Balance Gavish-Graves)} \NR
\stopalign
\stopformula
\stopsubsection

\startsubsection[title={Tiempo de ejecución y completitud}]
\placeformula
\startformula
t\_exec_{b} \ge B_{setup} + B_{picking} \cdot \left( \sum_{n \in \mathcal{N}} C_{n} \cdot Bn_{n b} \right) + B_{time} \cdot \left( \sum_{(i,j) \in \mathcal{E}} d_{ij} \cdot z_{i j b} \right) + \sum_{i \in \mathcal{V} \setminus \{v_1,v_f\}} S_{i} \cdot \left( \sum_{j:(i,j) \in \mathcal{E}} z_{i j b} \right)
\stopformula
\stopsubsection

\startsubsection[title={Asignación de lotes y secuenciación}]
\startformula
\startalign[n=3, align={left,left,left}]
    \NC \sum_{p \in \mathcal{P}} X_{p b} = 1; \NC \quad \forall b \in \mathcal{B} \NC \quad \text{(Asignación única por lote a un picker)} \NR
    \NC t\_exec_{b} - M \cdot (1 - X_{p b}) \le u_{p b}; \NC \quad \forall p \in \mathcal{P},\ b \in \mathcal{B} \NC \quad \text{(Ejecución acota procesamiento, cota inferior)} \NR
    \NC u_{p b} \le t\_exec_{b} + M \cdot (1 - X_{p b}); \NC \quad \forall p \in \mathcal{P},\ b \in \mathcal{B} \NC \quad \text{(Ejecución acota procesamiento, cota superior; fija } u_{pb} = t\_exec_{b} \text{ cuando } X_{pb}=1 \text{)} \NR
    \NC u_{p b} \le v_{p b}; \NC \quad \forall p \in \mathcal{P},\ b \in \mathcal{B} \NC \quad \text{(Procesamiento precede a completitud)} \NR
    \NC v_{p b} - M \cdot (1 - X_{p b}) \le v_{b}; \NC \quad \forall p \in \mathcal{P},\ b \in \mathcal{B} \NC \quad \text{(Completitud del lote vs picker)} \NR
\stopalign
\stopformula
\stopsubsection

\startsubsection[title={Secuenciamiento sin solapamiento}]
\placeformula
\startformula
\startalign[n=2, align={left,left}]
    \NC v_{p k} \ge v_{p b} + t\_exec_{k} - M_{seq} \cdot (3 - X_{p b} - X_{p k} - Z_{b k}); \NC \quad \forall p \in \mathcal{P},\ b,k \in \mathcal{B},\ b < k \NR
    \NC v_{p b} \ge v_{p k} + t\_exec_{b} - M_{seq} \cdot (2 - X_{p b} - X_{p k} + Z_{b k}); \NC \quad \forall p \in \mathcal{P},\ b,k \in \mathcal{B},\ b < k \NR
\stopalign
\stopformula

Donde $M_{seq}$ es una constante local de Big-M usada para linealizar el ordenamiento de lotes.
\stopsubsection

\startsubsection[title={Earliness y tardiness}]
\startformula
\startalign[n=3, align={left,left,left}]
    \NC v_{p b} - Dd_{n} - M \cdot (2 - X_{p b} - Bn_{n b}) \le ta_{n}; \NC \quad \forall p \in \mathcal{P},\ b \in \mathcal{B},\ n \in \mathcal{N} \NC \quad \text{(Definición de tardanza)} \NR
    \NC Dd_{n} - v_{p b} - M \cdot (2 - X_{p b} - Bn_{n b}) \le ea_{n}; \NC \quad \forall p \in \mathcal{P},\ b \in \mathcal{B},\ n \in \mathcal{N} \NC \quad \text{(Definición de adelanto)} \NR
\stopalign
\stopformula
\stopsubsection

\startsubsection[title={Dominio de variables}]
\startformula
\startalign[n=2, align={left,left}]
    \NC ta_{n}, ea_{n}, t\_exec_{b}, v_{b}, u_{p b}, v_{p b}, f_{i j b} \ge 0; \NC \quad \forall n \in \mathcal{N},\ b \in \mathcal{B},\ p \in \mathcal{P},\ (i,j) \in \mathcal{E} \NR
    \NC z_{i j b}, X_{p b}, Z_{b k} \in \{0,1\}; \NC \quad \forall p \in \mathcal{P},\ b,k \in \mathcal{B},\ (i,j) \in \mathcal{E} \NR
\stopalign
\stopformula
\stopsubsection

\stopsection

\stoptext